\section{Units, numeric primitives, and coordinate utilities}
The utility layer supplies the shared mathematics used by geometry, physics, plots and export. \class{Coordinate} stores Cartesian values and an optional weight; that weight can represent mass in one calculation and normal-force weighting in another. \class{WorldCoordinate} represents geographic location. \class{Quaternion} and \class{Transformation} encode orientation and rigid transforms, while inertia and bounding-box helpers support physical aggregation and geometry. These types do not remove the need to know the frame and units expected by the caller.

Interpolation, polynomial fitting, numerical statistics, array operations, noise generation and mathematical tolerance helpers are centralized in \code{util}. A linear interpolator over a thrust curve, a polynomial interpolation across aerodynamic regimes, and a pink-noise wind process solve different problems even though all are numerical helpers. Their boundary behavior, extrapolation rules and random-state ownership are part of the calling subsystem's contract.

The units package also does more than multiply by a display scale. Temperature conversions can include an offset. Caliber units depend on a rocket reference diameter, and percentage-of-length units depend on a relevant length. Fractional inch formatting and axis tick generation are presentation concerns. \class{FixedUnitGroup} gives custom quantities a fixed unit representation; it does not supply automatic physical dimension checking for arbitrary expressions. Values and comparators keep unit-aware display/sorting separate from internal SI calculations.

\sourceline{\source{core/src/main/java/info/openrocket/core/util/Coordinate.java}{Coordinate}; \source{core/src/main/java/info/openrocket/core/util/Quaternion.java}{Quaternion}; \source{core/src/main/java/info/openrocket/core/util/Transformation.java}{Transformation}; \source{core/src/main/java/info/openrocket/core/unit/UnitGroup.java}{UnitGroup}; \source{core/src/main/java/info/openrocket/core/unit/CaliberUnit.java}{caliber conversion}.}

\section{Change tracking, traversal, formatting, and object lifetimes}
Modification IDs identify changes that invalidate cached geometry or results. Listener lists and change-source helpers distribute notifications; invalidation helpers make disposed UI bindings stop acting on old models. Mutability guards distinguish objects still being assembled from completed data that should be read. Safety checks can detect inappropriate access, but they do not make a mutable document or engine safe for arbitrary concurrent writes.

Rocket visitors provide reusable traversals, including component and motor-mount listing and breadth/depth traversal strategies. They let consumers such as reports or selectors visit the component tree without embedding every traversal in a dialog. Active physical instance traversal remains a separate question from visiting each design node once.

The formatting package creates readable rocket/configuration descriptions and substitutes motor/configuration information into names. It is separate from expression evaluation and unit conversion. Configuration naming can depend on selected motors, while a persistent configuration ID remains the stable reference. Localization wraps translators, contextual keys and resource bundles so these labels can change with language without redefining stored identities.

File, JAR, reflection, platform, build-property and configuration helpers support startup and persistence. Their small size should not obscure their effect: resource lookup and class loading can determine whether a packaged application behaves like a development launch. The jar-in-jar bootstrap classes and platform helpers belong to this deployment boundary.

\section{Warnings, application logs, and network services}
Simulation warnings/errors are structured model output with types and associated context. They can survive in results and appear in reports. Application logging is a different channel: Swing logging classes integrate log buffers, filtering, levels and diagnostic displays with the logging backend. A debug message about a repaint is not a flight warning; a flight warning should not disappear merely because a console logging threshold changes.

Communication services retrieve welcome/release information and update metadata, and support bug-report workflows. Connection abstractions separate network access from response interpretation. MIT-specific update selection and replacement are described in Chapter~\ref{ch:build}; inherited update classes also remain in the tree. The existence of a network helper does not mean every simulation uses the network. Normal physics calculations operate on the already resolved design, motor and environment inputs.

ThrustCurve communication has its own request/response and curve decoding classes. A motor download must become a validated local motor record before it can participate in a configuration. Network metadata, local catalog entries and a runtime motor's ignition clock have separate lifetimes and failure modes.

\sourceline{\source{core/src/main/java/info/openrocket/core/logging/WarningSet.java}{structured warnings}; \source{core/src/main/java/info/openrocket/core/communication/ConnectionSource.java}{connection abstraction}; \source{core/src/main/java/info/openrocket/core/communication/WelcomeInfoRetriever.java}{welcome information}; \source{core/src/main/java/info/openrocket/core/formatting/RocketDescriptorImpl.java}{rocket description formatting}.}

\section{Inherited extension examples and developer utilities}
The extension example packages illustrate air starts, CSV saving, damping-moment changes, simulation printing, roll control and stopping conditions. Providers expose examples to discovery; runtime listeners implement their behavior. These examples are useful API references, but merely finding a class in the tree does not establish that a particular simulation has enabled it. The same distinction applies to Swing extension configurators: a panel edits configuration, while the core extension installs execution behavior.

The scripting support includes script-engine discovery and a GraalJS factory in addition to the simulation extension itself. Language availability and the packaged runtime therefore affect script execution. The separately bundled expression parser is for arithmetic/array expressions and does not replace the script-engine API or the MIT Java-board runtime.

Developer utility packages contain localization generators/reports, motor-correlation tools, optimizer exercises, preset conversion/editing/serialization utilities, and small desktop support tools. They are not all normal startup services. Build scripts, resources, tests and documentation complete the codebase around the runtime Java packages; Appendix~\ref{app:packages} provides an explicit coverage inventory so these less visible areas are not omitted from the architecture map.
